\clearpage
\section{Detailed derivations from the course}
\label{app:derivations}
% EDITORIAL ADDITION: navigation or conceptual bridge
These longer calculations support the conceptual development in the main text.
Their algebra, assumptions, and interpretations are retained from the day-by-day notes.
% Source: day3.tex:133; unit d3.0.3
\subsection{Discrete Fourier modes and the long-wavelength decay law}
\label{deriv:discrete-modes}

% ====================

Consider one cosine mode,
\[
    P_n(x,t) = A_n\cos(q_nx),
\]
with
\[
    q_n = \frac{2\pi n}{L}.
\]
Substituting this into the random-walk update rule gives
\begin{align*}
    P_n(x,t+\Dt) &= A_n \cos(q_nx) + k\Dt A_n \cos \left( q_n(x-a) \right)
    \nonumber\\
    &\quad + k\Dt A_n \cos \left( q_n(x+a) \right) - 2k\Dt A_n \cos(q_nx).
\end{align*}
Using
\[
    \cos (\alpha + \beta) + \cos (\alpha - \beta) = 2 \cos \alpha \cos \beta,
\]
we obtain
\begin{align*}
    P_n(x,t+\Dt) &= A_n\cos(q_nx) + 2k\Dt A_n\cos(q_nx)\cos(q_na)
    \nonumber\\
    &\qquad - 2k\Dt A_n\cos(q_nx)
    \\
    &= A_n\cos(q_nx) \left[ 1+2k\Dt\left(\cos(q_na)-1\right) \right].
\end{align*}
Since
\[
    P_n(x,t) = A_n\cos(q_nx),
\]
we have
\[
    \boxed{ P_n(x,t+\Dt) = r_n P_n(x,t), }
\]
where
\[
    \boxed{ r_n = 1 - 2k\Dt \left[ 1-\cos(q_na) \right]. }
\]
The spatial shape $\cos(q_nx)$ has not changed. 
Only its amplitude has been multiplied by the number $r_n$.

Therefore,
\[
    \boxed{ \cos(q_nx) \text{ is an eigenfunction of the random-walk update operator.} }
\]
The corresponding eigenvalue for one time step is $r_n$.
% 
% 
% ====================
\medskip
\\
\textbf{Physical meaning of the Fourier modes}
% ====================

Each value of $n$ corresponds to a different spatial wavelength,
\[
    \lambda_n = \frac{2\pi}{q_n} = \frac{L}{n}.
\]
Small $n$ corresponds to long-wavelength, slowly varying structures. 
Large $n$ corresponds to short-wavelength, rapidly varying structures. 
Diffusion suppresses the short-wavelength modes more rapidly than the long-wavelength modes. 
This is why an initially sharp distribution becomes smooth with time.
% 
% 
% ====================
\medskip
\\
\textbf{Small-angle approximation and the diffusion constant}
% ====================

For wavelengths much larger than the lattice spacing,
\[
    q_na \ll 1.
\]
Using the Taylor expansion
\[
    \cos\theta = 1 - \frac{\theta^2}{2} + \mathcal{O}(\theta^4),
\]
we obtain
\[
    \cos(q_na) \approx 1 - \frac{q_n^2a^2}{2}.
\]
Therefore,
\begin{align*}
    r_n &= 1 - 2k\Dt \left[ 1-\cos(q_na) \right]
    \\
    &\approx 1 - 2k\Dt \left( \frac{q_n^2a^2}{2} \right)
    \\
    &= 1-ka^2q_n^2\Dt.
\end{align*}
Define the diffusion constant
\[
    \boxed{ D = ka^2. }
\]
Then
\[
    \boxed{ r_n \approx 1-Dq_n^2\Dt. }
\]
Hence
\[
    P_n(x,t + \Dt) = P_n(x,t) \left( 1 - Dq_n^2 \Dt \right).
\]
% 
% 
% ====================
\medskip
\\
\textbf{Evolution over many time steps}
% ====================

After $N$ time steps,
% 
\[
    P_n(x,t+N\Dt) = r_n^N P_n(x,t).
\]
% 
In the long-wavelength approximation,
% 
\[
    r_n^N = \left( 1-Dq_n^2\Dt \right)^N.
\]
% 
Let the total elapsed time be
% 
\[
    \tau = N\Dt.
\]
% 
Therefore,
% 
\[
    N = \frac{\tau}{\Dt},
\]
% 
and
% 
\[
    r_n^N = \left( 1-Dq_n^2\Dt \right)^{\tau/\Dt}.
\]
% 
Using
% 
\[
    \lim_{\epsilon\rightarrow0} (1+\alpha\epsilon)^{1/\epsilon} = e^\alpha,
\]
% 
we obtain
% 
\[
    \boxed{ r_n^N \longrightarrow e^{-Dq_n^2\tau}. }
\]
% 
Thus a Fourier mode evolves as
% 
\[
    \boxed{ P_n(x,t) = A_n e^{-Dq_n^2t} \cos(q_nx). }
\]
% 
The relaxation time of mode $n$ is therefore
% 
\[
    \tau_n = \frac{1}{Dq_n^2}.
\]
% 
Since $q_n$ increases with $n$, high-frequency modes decay much faster than low-frequency modes.


% ====================

% Source: day3.tex:691; unit d3.6.4
\subsection{The continuum limit: symmetry, Taylor expansion, and assumptions}
\label{deriv:continuum}

% ------------------------------------------------------------

Start from the discrete random-walk update rule,

\[
    P(x,t+\Delta t) = P(x,t) + k\Delta t\,P(x-a,t) + k\Delta t\,P(x+a,t) - 2k\Delta t\,P(x,t).
\]

Equivalently,

\[
    P(x,t+\Delta t) = P(x,t) + k\Delta t\left[P(x-a,t)+P(x+a,t)-2P(x,t)\right].
\]
Here $k$ is the hopping rate in each direction and $a$ is the lattice spacing.

We now assume that $P(x,t)$ varies smoothly over distances comparable to $a$.
This allows us to Taylor-expand the shifted probability distributions around $x$.
For the right-shifted position,
% % 
\[
    P(x+a,t) = P(x,t) + a\frac{\partial P}{\partial x} + \frac{a^2}{2}\frac{\partial^2P}{\partial x^2} + \frac{a^3}{3!}\frac{\partial^3P}{\partial x^3} + \cdots.
\]
% % 
For the left-shifted position,
% % 
\[
    P(x-a,t) = P(x,t) - a\frac{\partial P}{\partial x} + \frac{a^2}{2}\frac{\partial^2P}{\partial x^2} - \frac{a^3}{3!}\frac{\partial^3P}{\partial x^3} + \cdots.
\]
% 
Adding the two expansions causes all odd-order derivatives to cancel,
% 
\[
    P(x+a,t)+P(x-a,t) = 2P(x,t) + a^2\frac{\partial^2P}{\partial x^2} + \mathcal{O}(a^4).
\]
% 
Therefore,
% 
\[
    P(x-a,t)+P(x+a,t)-2P(x,t) = a^2\frac{\partial^2P}{\partial x^2} + \mathcal{O}(a^4).
\]
% 
Substituting this result into the update rule gives
% 
\[
    P(x,t+\Delta t) = P(x,t) + k\Delta t\,a^2\frac{\partial^2P}{\partial x^2} + \mathcal{O}(a^4).
\]
% 
Now Taylor-expand the left-hand side in time,
% 
\[
    P(x,t+\Delta t) = P(x,t) + \Delta t\frac{\partial P}{\partial t} + \mathcal{O}(\Delta t^2).
\]
% 
Substituting this gives
% 
\[
    P(x,t) + \Delta t\frac{\partial P}{\partial t} = P(x,t) + k\Delta t\,a^2\frac{\partial^2P}{\partial x^2}.
\]
% 
Canceling $P(x,t)$ from both sides gives
% 
\[
    \Delta t\frac{\partial P}{\partial t} = k\Delta t\,a^2\frac{\partial^2P}{\partial x^2}.
\]
% 
Canceling $\Delta t$ gives
% 
\[
    \frac{\partial P}{\partial t} = ka^2\frac{\partial^2P}{\partial x^2}.
\]
% 
Defining
% 
\[
    \boxed{D=ka^2,}
\]
% 
we recover
% 
\[
    \boxed{\frac{\partial P}{\partial t} = D\frac{\partial^2P}{\partial x^2}.}
\]
% 
Thus the continuum diffusion equation is the long-wavelength, smooth-space limit of the discrete random walk.
% 
% 
% ------------------------------------------------------------
\medskip
\\
\textbf{Why the first derivative disappears}
% ------------------------------------------------------------

The cancellation of the first spatial derivative is a consequence of the symmetry of the random walk.
The walker moves left and right with the same rate.
The term proportional to $+\partial P/\partial x$ from one direction therefore cancels the term proportional to $-\partial P/\partial x$ from the other.
This leaves the second derivative as the leading spatial contribution,

\[
    \frac{\partial^2P}{\partial x^2}.
\]

If the walker instead had a systematic preference for one direction, the first derivatives would not cancel completely.
The continuum equation would then contain a drift term in addition to diffusion.
% 
% 
The continuum-limit derivation assumes that:
\begin{itemize}
    \item the walk is unbiased, so left and right hopping rates are equal;
    \item the hopping rate $k$ is independent of position and time;
    \item only nearest-neighbor jumps of length $a$ occur;
    \item the process is Markovian, so the next transition depends only on the current state;
    \item $P(x,t)$ varies smoothly on the length scale $a$, allowing the spatial Taylor expansion;
    \item $\Delta t$ is sufficiently small that the temporal Taylor expansion is valid;
    \item higher-order terms in $a$ and $\Delta t$ can be neglected;
    \item the continuum limit is taken while the combination $ka^2=D$ remains finite.
\end{itemize}


% ------------------------------------------------------------

% Source: day3.tex:403; unit d3.3
\subsection{Solving diffusion by separation of variables}
\label{deriv:separation}


% ====================

We now recover the Fourier modes directly from the PDE.

Start from
% 
\[
    \frac{\partial P}{\partial t} = D\frac{\partial^2P}{\partial x^2}.
\]
% 
Assume a separable solution
% 
\[
    P(x,t) = X(x) T(t).
\]
% 
Substitution gives
% 
\[
    X(x)\frac{\dd T}{\dd t} = DT(t)\frac{\dd^2X}{\dd x^2}.
\]
% 
Divide by $DXT$:
% 
\[
    \frac{1}{DT} \frac{\dd T}{\dd t} = \frac{1}{X} \frac{\dd^2X}{\dd x^2}.
\]
% 
The left-hand side depends only on $t$, while the right-hand side depends only on $x$. 
They can therefore be equal for every $x$ and $t$ only if both are equal to the same constant.
For diffusive modes, write this constant as $-q^2$:
% 
\[
    \frac{1}{DT} \frac{\dd T}{\dd t} = \frac{1}{X} \frac{\dd^2X}{\dd x^2} = -q^2.
\]
% 
This gives two ordinary differential equations:
% 
\begin{align*}
    \frac{\dd T}{\dd t} &= -Dq^2T,
    \\
    \frac{\dd^2X}{\dd x^2} &= -q^2X.
\end{align*}
% 
Their solutions are
% 
\[
    T(t) = Ce^{-Dq^2t},
\]
% 
and
% 
\[
    X(x) = A\cos(qx) + B\sin(qx).
\]
% 
Therefore,
% 
\[
    P(x,t) = \left[ A\cos(qx)+B\sin(qx) \right] e^{-Dq^2t}.
\]

Periodic boundary conditions require
% 
\[
    q = q_n = \frac{2\pi n}{L}.
\]
% 
Thus
% 
\[
    \boxed{ P_n(x,t) = \left[ A_n\cos(q_nx) + B_n\sin(q_nx) \right] e^{-Dq_n^2t}. }
\]
% 
% 
% ====================
\medskip
\\
\textbf{Why this solves a general initial condition}
% ====================

The usefulness of Fourier analysis is not merely that a cosine happens to solve the equation.
A sufficiently well-behaved periodic initial probability distribution can be decomposed into Fourier modes,

\[
    P(x,0) = A_0 + \sum_{n=1}^{\infty} \left[ A_n\cos(q_nx) + B_n\sin(q_nx) \right].
\]
% 
Because the diffusion equation is linear, each Fourier mode evolve independently:
% 
\[
    P(x,t) = A_0 + \sum_{n=1}^{\infty} \left[ A_n\cos(q_nx) + B_n\sin(q_nx) \right] e^{-Dq_n^2t}.
\]
% 
Thus the original PDE has effectively been converted into independent equations for the amplitudes of the Fourier modes.
% 
For each mode,
% 
\[
    \frac{\dd A_n}{\dd t} = -Dq_n^2A_n,
\]
% 
which is just an ordinary differential equation.
% 
This is the central simplification:
% 
\[
    \boxed{ \text{PDE in real space} \quad\longrightarrow\quad \text{independent ODEs for Fourier modes}. }
\]


% ====================

% Source: day5.tex:211; unit d5.2.2
\subsection{Straight-line least squares: deriving the normal equations}
\label{deriv:fit}

% ------------------------------------------------------------

At the minimum of the error function,

\[
    \frac{\partial E}{\partial a}=0
\]
% 
and
% 
\[
    \frac{\partial E}{\partial b}=0.
\]
% 
Differentiating with respect to $a$ gives
% 
\[
    \frac{\partial E}{\partial a}=\sum_{i=1}^{N}-2x_i\left[y_i-(ax_i+b)\right]=0.
\]
% 
Therefore,
% 
\[
    a\sum_{i=1}^{N}x_i^2+b\sum_{i=1}^{N}x_i=\sum_{i=1}^{N}x_iy_i.
\]
% 
Differentiating with respect to $b$ gives
% 
\[
    \frac{\partial E}{\partial b}=\sum_{i=1}^{N}-2\left[y_i-(ax_i+b)\right]=0.
\]
% 
Therefore,
% 
\[
    a\sum_{i=1}^{N}x_i+Nb=\sum_{i=1}^{N}y_i.
\]
% 
The fitting problem has now become a system of two linear equations for the two unknown parameters $a$ and $b$.


% ------------------------------------------------------------

% Source: day5.tex:254; unit d5.2.3
\subsubsection{Matrix representation and explicit fitting parameters}
\label{deriv:fit-matrix}

% ------------------------------------------------------------

The two equations can be written as
% 
\[
    \begin{pmatrix}
        \sum_i x_i^2 & \sum_i x_i \\
        \sum_i x_i & N
    \end{pmatrix}
    \begin{pmatrix}
        a \\
        b
    \end{pmatrix}
    =
    \begin{pmatrix}
        \sum_i x_i y_i \\
        \sum_i y_i
    \end{pmatrix}.
\]
% 
A general system of linear equations has the form
% 
\[
    A\mathbf{x}=\mathbf{c}.
\]
% 
Formally, if $A$ is invertible,
% 
\[
    \mathbf{x}=A^{-1}\mathbf{c}.
\]
% 
For the straight-line fit, define
% 
\[
    \Delta=N\sum_i x_i^2-\left(\sum_i x_i\right)^2.
\]
% 
The slope is
% 
\[
    \boxed{a=\frac{N\sum_i x_i y_i-\left(\sum_i x_i\right)\left(\sum_i y_i\right)}{\Delta}.}
\]
% 
The intercept is
% 
\[
    \boxed{b=\frac{\left(\sum_i x_i^2\right)\left(\sum_i y_i\right)-\left(\sum_i x_i\right)\left(\sum_i x_i y_i\right)}{\Delta}.}
\]
% 
% 
% ============================================================

% Source: day5.tex:558; unit d5.5.1
\subsection{Maximum entropy with normalization and a mean constraint}
\label{deriv:maxent}

% ------------------------------------------------------------

Suppose that the states are labeled
% 
\[
    i=1,\ldots,n.
\]
% 
The probabilities must satisfy the normalization constraint
% 
\[
    \sum_{i=1}^{n}p_i=1.
\]
% 
As an example, suppose that the mean state is also constrained to be
% 
\[
    \sum_{i=1}^{n}ip_i=4.5.
\]
% 
We want to find the distribution $\{p_i\}$ that maximizes
% 
\[
    S=-\sum_{i=1}^{n}p_i\ln p_i
\]
% 
subject to these constraints.
% 
Introduce Lagrange multipliers $\alpha$ and $\beta$ and define
% 
\[
    \mathcal{L}=-\sum_{i=1}^{n}p_i\ln p_i-\alpha\left(\sum_{i=1}^{n}p_i-1\right)-\beta\left(\sum_{i=1}^{n}ip_i-4.5\right).
\]
% 
The stationary condition for every $p_i$ is
% 
\[
    \frac{\partial\mathcal{L}}{\partial p_i}=0.
\]
% 
This gives
% 
\[
    -\ln p_i-1-\alpha-\beta i=0.
\]
% 
Therefore,
% 
\[
    p_i=e^{-1-\alpha-\beta i}.
\]
% 
Equivalently,
% 
\[
    p_i\propto e^{-\beta i}.
\]
% 
The remaining constants are determined by the normalization and mean-value constraints.


% ============================================================

